% from ega3-1-fr.tex
\oldpage[III]{135}
déduit que si $f:X\xrightarrow{f'}Y'\xrightarrow{g}Y$ est la factorisation
de Stein de $f$ (\hyperref[III.4.3.3-fr]{4.3.3}), $g^{-1}(y)$ se réduit à
un seul point $y'$~; en outre, on a
$\mathcal O_y\subset\mathcal O_{y'}=(f_*(\mathcal O_X))_y
\subset\mathrm R(X)$,
et comme $\mathcal O_{y'}$ est fini sur $\mathcal O_y$, (et \emph{a
fortiori} sur $\mathrm R(Y)$), il est contenu dans $\mathrm R(Y)$ en vertu
de l'hypothèse~; comme $y$ est normal, on a nécessairement
$\mathcal O_{y'}=\mathcal O_y$~; on en conclut que $g$ est un isomorphisme
local au point $y'$ (\textbf{I}, 6.5.4), ce qui achève de prouver la première partie
du corollaire. La seconde résulte de la première, car l'hypothèse
supplémentaire entraîne que $g$ est bijective et un isomorphisme local au
voisinage de tout point de $Y'$, donc un isomorphisme.
\end{proof}

Le fait que (\hyperref[III.4.3.7-fr]{4.3.7}) soit établi dans le cadre des
schémas permet des applications telles que la suivante~:

\begin{proposition}[4.3.13]
\label{III.4.3.13-fr}
Soient $A$ un anneau local noethérien unibranche, $\mathfrak a$ un idéal de
définition de $A$, $A_0=A/\mathfrak a$,
$S=\operatorname{gr}_{\mathfrak a}(A)$ l'anneau gradué associé à $A$ pour
la filtration $\mathfrak a$-préadique~; $S$ est une $A_0$-algèbre graduée
engendrée par $S_1$, $S_1$ étant un $A_0$-module de type fini. Alors
$\operatorname{Proj}(S)$ est un $A_0$-schéma connexe.
\end{proposition}

\begin{proof}
Soit $\mathfrak m$ l'idéal maximal de $A$~; $Y=\operatorname{Spec}(A)$ est
un schéma intègre dont le point $y$ correspondant à $\mathfrak m$ est
l'unique point fermé. Par hypothèse, on a
$\mathfrak m^p\subset\mathfrak a\subset\mathfrak m$ pour un entier $p$,
donc $V(\mathfrak a)=\{\mathfrak m\}$. Soit
$S'=\bigoplus_{n\geq0}\mathfrak a^n$, et soit
$X=\operatorname{Proj}(S')$, qui est le $Y$-schéma obtenu en faisant éclater
l'idéal $\mathfrak a$~; $X$ est intègre et le morphisme structural
$f:X\to Y$ est birationnel (\textbf{II}, 8.1.4) et évidemment projectif. Par suite,
(\hyperref[III.4.3.7-fr]{4.3.7}) est applicable et montre que $f^{-1}(y)$
est connexe~; mais l'espace $f^{-1}(y)$ est sous-jacent à
$\operatorname{Proj}(S'\otimes_AA_0)$ (\textbf{I}, 3.6.1 et \textbf{II}, 2.8.10)~; comme
$S'\otimes_AA_0=S$ par définition, la proposition est démontrée.
\end{proof}


\subsection{Le «~main theorem~» de Zariski.}
\label{subsection:III.4.4-fr}

\begin{proposition}[4.4.1]
\label{III.4.4.1-fr}
Soient $Y$ un préschéma localement noethérien, $f:X\to Y$ un morphisme
propre. Soit $X'$ l'ensemble des points $x\in X$ qui sont isolés dans leur
fibre $f^{-1}(f(x))$. Alors l'ensemble $X'$ est ouvert dans $X$, et si
$f=g\circ f'$ est la factorisation de Stein de $f$
\textup{(\hyperref[III.4.3.3-fr]{4.3.3})}, la restriction de $f'$ à $X'$ est un
isomorphisme de $X'$ sur un sous-préschéma induit sur un ouvert $U$ de $Y'$,
et on a $X'=f'^{-1}(U)$.
\end{proposition}

\begin{proof}
Comme $g^{-1}(f(x))$ est fini et discret
(\hyperref[III.4.3.3-fr]{4.3.3} et \textbf{II}, 6.1.7), pour que $x$ soit isolé dans
$f^{-1}(f(x))$, il faut et il suffit qu'il le soit dans
$f'^{-1}(f'(x))$~; on peut donc se borner au cas où $f'=f$, donc
$f_*(\mathcal O_X)=\mathcal O_Y$. Alors, si $x\in X'$,
$f^{-1}(f(x))$, qui est connexe (\hyperref[III.4.3.2-fr]{4.3.2}) est
nécessairement réduit au point $x$. Comme $f$ est fermé, pour tout voisinage
ouvert $V$ de $x$ dans $X$, $f(X\setmin V)$ est fermé dans $Y$ et ne contient pas
$y=f(x)$, puisque $f^{-1}(y)=\{x\}$~; si $U$ est le complémentaire de
$f(X\setmin V)$ dans $Y$, on a $f^{-1}(U)\subset V$, et on en conclut que les
images réciproques par $f$ d'un système fondamental de voisinages ouverts de
$y$ forment un système fondamental de voisinages ouverts de $x$. L'hypothèse
$f_*(\mathcal O_X)=\mathcal O_Y$ et la définition de l'image directe d'un
faisceau (\textbf{0}\textsubscript{I}, 3.4.1 et 4.2.1) entraînent alors que, si
$f=(\psi,\theta)$, l'homomorphisme
$\theta_y^\sharp:\mathcal O_y\to\mathcal O_x$ est un \emph{isomorphisme}. On en conclut
qu'il existe un voisinage ouvert $V$ de $x$ et un voisinage ouvert $U$ de
$y$ tels que la restriction de $f$ à $V$ soit un isomorphisme de $V$ sur
$U$ (\textbf{I}, 6.5.4)~; en outre, d'après ce qu'on vient de voir,
